\section{Thiết kế cơ chế gọi RAG thông qua Custom Action}

Trong hệ thống, RAG không được tích hợp trực tiếp vào Rasa Core mà được gọi thông qua Custom Action. Cách thiết kế này giúp tách biệt rõ giữa phần điều phối hội thoại của Rasa và phần truy xuất tri thức của RAG.

\subsection*{Vai trò của RAG trong chatbot tuyển sinh}

RAG được sử dụng để xử lý các câu hỏi:
\begin{itemize}
    \item Có nội dung dài hoặc phức tạp.
    \item Cần truy xuất từ tài liệu tuyển sinh.
    \item Có khả năng thay đổi theo từng năm.
    \item Không phù hợp để viết thành câu trả lời tĩnh.
    \item Cần trích dẫn nguồn từ tài liệu.
\end{itemize}

Ví dụ các câu hỏi nên gọi RAG:
\begin{itemize}
    \item Quy định xét tuyển thẳng năm nay như thế nào?
    \item Điều kiện đăng ký xét tuyển bằng chứng chỉ quốc tế là gì?
    \item Hãy giải thích chi tiết đề án tuyển sinh của trường.
    \item Ngành công nghệ thông tin khác gì ngành an toàn thông tin?
\end{itemize}

\subsection*{Luồng gọi RAG}

\begin{verbatim}
Người dùng
   ↓
Frontend Web
   ↓
Backend API
   ↓
Rasa Server
   ↓
Custom Action: action_call_rag
   ↓
RAG Service
   ↓
Vector Store / Document Store / Knowledge Graph
   ↓
LLM sinh câu trả lời
   ↓
Custom Action
   ↓
Rasa Server
   ↓
Backend API
   ↓
Frontend Web
   ↓
Người dùng
\end{verbatim}

\subsection*{Thiết kế request từ Custom Action sang RAG Service}

\begin{verbatim}
{
  "query": "Quy định xét tuyển thẳng năm nay như thế nào?",
  "conversation_id": "conv_001",
  "user_id": "user_001",
  "slots": {
    "major": "công nghệ thông tin",
    "year": "2025"
  },
  "metadata": {
    "source": "rasa",
    "language": "vi"
  }
}
\end{verbatim}

\begin{table}[H]
\centering
\begin{tabular}{|l|p{9cm}|}
\hline
\textbf{Trường} & \textbf{Ý nghĩa} \\
\hline
query & Câu hỏi gốc của người dùng \\
\hline
conversation\_id & Mã hội thoại \\
\hline
user\_id & Mã người dùng \\
\hline
slots & Thông tin ngữ cảnh Rasa đã trích xuất \\
\hline
metadata & Thông tin bổ sung như nguồn gọi, ngôn ngữ \\
\hline
\end{tabular}
\end{table}

\subsection*{Thiết kế response từ RAG Service}

\begin{verbatim}
{
  "answer": "Theo đề án tuyển sinh, thí sinh đạt giải học sinh giỏi cấp quốc gia có thể được xét tuyển thẳng theo quy định của Bộ Giáo dục và Đào tạo...",
  "sources": [
    {
      "title": "Đề án tuyển sinh 2025",
      "page": 12,
      "url": "/documents/de-an-tuyen-sinh-2025.pdf",
      "score": 0.87
    }
  ],
  "confidence": 0.82,
  "answer_type": "rag"
}
\end{verbatim}

\subsection*{Ví dụ action\_call\_rag}

\begin{verbatim}
class ActionCallRAG(Action):

    def name(self) -> Text:
        return "action_call_rag"

    def run(
        self,
        dispatcher: CollectingDispatcher,
        tracker: Tracker,
        domain: Dict[Text, Any]
    ) -> List[Dict[Text, Any]]:

        user_query = tracker.latest_message.get("text")
        sender_id = tracker.sender_id

        major = tracker.get_slot("major")
        year = tracker.get_slot("year")

        payload = {
            "query": user_query,
            "conversation_id": sender_id,
            "slots": {
                "major": major,
                "year": year
            },
            "metadata": {
                "source": "rasa",
                "language": "vi"
            }
        }

        try:
            response = requests.post(
                f"{RAG_SERVICE_URL}/rag/query",
                json=payload,
                timeout=15
            )

            if response.status_code != 200:
                dispatcher.utter_message(
                    text="Mình chưa truy xuất được tài liệu phù hợp. Bạn vui lòng thử lại sau."
                )
                return [SlotSet("last_answer_type", "fallback")]

            result = response.json()

            answer = result.get("answer")
            sources = result.get("sources", [])
            confidence = result.get("confidence", 0)

            if not answer or confidence < 0.5:
                dispatcher.utter_message(
                    text=(
                        "Mình chưa tìm thấy thông tin đủ chắc chắn trong tài liệu. "
                        "Bạn có thể hỏi lại cụ thể hơn, ví dụ: ngành nào, năm nào hoặc phương thức xét tuyển nào."
                    )
                )
                return [SlotSet("last_answer_type", "fallback")]

            message = answer

            if sources:
                message += "\n\nNguồn tham khảo:"
                for index, source in enumerate(sources[:3], start=1):
                    title = source.get("title", "Tài liệu tuyển sinh")
                    page = source.get("page")
                    if page:
                        message += f"\n{index}. {title}, trang {page}"
                    else:
                        message += f"\n{index}. {title}"

            dispatcher.utter_message(text=message)

            return [
                SlotSet("rag_query", user_query),
                SlotSet("last_answer_type", "rag")
            ]

        except Exception:
            dispatcher.utter_message(
                text=(
                    "Hệ thống truy xuất tài liệu hiện chưa phản hồi. "
                    "Bạn có thể thử lại sau hoặc hỏi các thông tin phổ biến như ngành học, điểm chuẩn, học phí."
                )
            )

            return [SlotSet("last_answer_type", "fallback")]
\end{verbatim}

\subsection*{Ưu điểm của thiết kế gọi RAG qua Custom Action}

Thiết kế này có các ưu điểm sau:
\begin{itemize}
    \item Rasa chỉ tập trung vào hiểu intent và điều phối hội thoại.
    \item RAG Service hoạt động độc lập, dễ nâng cấp mà không ảnh hưởng đến Rasa.
    \item Có thể thay đổi mô hình embedding, vector database hoặc LLM mà không cần thay đổi logic hội thoại.
    \item Câu trả lời có thể kèm nguồn tài liệu, giúp tăng độ tin cậy.
    \item Hệ thống có khả năng mở rộng khi số lượng tài liệu tuyển sinh tăng lên.
\end{itemize}
